\cleardoublepage % memastikan bab baru dimulai di halaman ganjil (kanan)
\chapter{STUDI LITERATUR}
\label{chap:studi-literatur}

Bab ini merupakan hasil tahap \textit{Empathize} dari kerangka \textit{Design
	Thinking} yang dipakai pada tugas akhir ini (lihat Subbab 1.5). Pembahasan
disusun dari yang paling umum ke yang paling spesifik terhadap persoalan
tugas akhir ini: tren elektrifikasi kendaraan secara luas (Subbab
\ref{sec:elektrifikasi-kendaraan}), tantangan operasional bus listrik
(Subbab \ref{sec:bus-listrik}), model konsumsi energi baterainya (Subbab
\ref{sec:baterai-bus-listrik}), lalu perencanaan infrastruktur pengisian daya
sebagai persoalan optimasi maupun simulasi (Subbab
\ref{sec:perencanaan-infrastruktur}), dan ditutup dengan posisi tugas akhir
ini terhadap peta riset yang ada (Subbab \ref{sec:posisi-penelitian}).

\section{Elektrifikasi Kendaraan: Konteks dan Tren}
\label{sec:elektrifikasi-kendaraan}

Peralihan dari kendaraan berbahan bakar fosil ke kendaraan listrik terus
berlangsung cepat. \textit{Global EV Outlook 2026} yang diterbitkan
International Energy Agency (IEA) mencatat penjualan mobil listrik dunia
tumbuh 20\% pada 2025, melampaui 20 juta unit dan menyentuh seperempat dari
seluruh penjualan mobil baru secara global \parencite{iea2026}. Di kawasan
Asia Tenggara, penjualan mobil listrik meningkat lebih dari dua kali lipat
dibanding tahun sebelumnya; Indonesia menjadi salah satu penggeraknya, dengan
penjualan mobil listrik mencapai 15\% dari total penjualan mobil baru pada
2025 dan tumbuh hampir 90\% secara tahunan pada kuartal pertama 2026,
didorong turunnya premium harga mobil listrik dibanding kendaraan konvensional
dari lebih 50\% pada 2024 menjadi sekitar 40\% pada 2025 \parencite{iea2026}.

%TODO-CITE: jika ada data/laporan resmi Kemenhub atau Dishub DIY soal
% roadmap elektrifikasi transportasi publik Indonesia secara spesifik,
% tambahkan sitasinya di sini untuk memperkuat konteks lokal (bukan cuma
% data global IEA).
Tren ini merambat ke sektor transportasi publik. Penjualan bus listrik
global tumbuh 12\% pada 2025, dengan jangkauan rata-rata model bus listrik
yang tersedia di pasar kini mencapai sekitar 360 km sekali pengisian penuh
-- naik 15\% dibanding 2020, dan sudah cukup untuk menutupi kebutuhan
operasional harian bus kota yang umumnya menempuh 150--300 km per hari
\parencite{iea2026}. Perkembangan ini turut ditopang oleh perluasan
infrastruktur \textit{depot charging} sebagai salah satu penopang utama
kelayakan operasional dan skalabilitas armada bus listrik \parencite{iea2026}.

Namun peralihan pada level armada -- bukan kendaraan individual -- membawa
persoalan perencanaan tersendiri yang tidak muncul pada kendaraan pribadi.
\textcite{perumal2022} memetakan proses
perencanaan bus listrik ke dalam tiga tingkatan yang berurutan: strategis
(investasi armada dan infrastruktur pengisian daya, termasuk penempatan
lokasinya), taktis (penjadwalan kendaraan terhadap trip yang sudah
ditetapkan), dan operasional (penjadwalan pengisian daya harian). Pada
tingkat strategis, mereka mencatat setidaknya empat teknologi pengisian daya
yang lazim dipakai -- \textit{plug-in} lambat di depot, \textit{plug-in}
cepat atau pantograf di terminal/halte, jalur kontak/induksi nirkabel
selama bus berjalan, dan pertukaran baterai (\textit{battery swapping}) --
yang masing-masing punya karakteristik biaya pemasangan dan daya pengisian
yang sangat berbeda. Pemetaan tingkatan perencanaan inilah yang menjadi
acuan tugas akhir ini dalam membatasi cakupan pembahasan pada tingkat
strategis (lokasi dan kapasitas SPKLU) dan operasional (mekanisme
pengisian daya), sebagaimana dijelaskan pada subbab berikut. SPKLU yang
dimaksud di sini adalah Stasiun Pengisian Kendaraan Listrik Umum, yaitu
infrastruktur yang menyediakan satu atau lebih \textit{gun} (konektor
pengisian) tempat bus mengisi ulang baterainya.

\section{Bus Listrik: Karakteristik Operasional dan Tantangannya}
\label{sec:bus-listrik}

Berbeda dari mobil listrik pribadi yang jadwal pengisian dayanya fleksibel,
bus listrik beroperasi dalam jadwal rit yang ketat sepanjang hari dan tetap
harus melayani penumpang selama pengisian daya berlangsung -- kecuali
ditarik sementara dari layanan. Perencanaan armada bus listrik karenanya
memerlukan pertimbangan yang tidak berlaku untuk kendaraan pribadi:
jangkauan yang cukup untuk menyelesaikan rit dalam sehari, serta waktu dan
lokasi pengisian daya bus tanpa mengganggu jadwal.

Dua pendekatan pengisian daya umum digunakan: \textit{depot charging}
(pengisian di garasi, biasanya di luar jam operasional) dan
\textit{opportunity charging} (pengisian singkat di sela jam operasional,
misalnya saat bus berhenti di ujung rute). \textcite{jefferiesgohlich2020} mencatat bahwa penggantian bus diesel dengan
bus listrik satu-banding-satu, tanpa penyesuaian jadwal rit, kerap tidak
dapat diterapkan pada praktiknya -- keterbatasan jangkauan dan waktu
pengisian daya memaksa perubahan pada jumlah armada maupun pola rit yang
sudah berjalan.

Persoalan penempatan infrastruktur pengisian daya menjadi lebih rumit ketika
elektrifikasi diterapkan pada skala jaringan, bukan satu rute tunggal.
\textcite{rogge2018} menunjukkan bahwa ukuran armada dan
kebutuhan kapasitas infrastruktur saling memengaruhi -- menambah unit
\textit{charging} di satu lokasi bisa mengurangi kebutuhan penambahan
armada, dan sebaliknya -- sehingga keduanya perlu dipertimbangkan bersamaan.
\textcite{he2023} membawa keterkaitan ini satu langkah
lebih jauh dengan menggabungkan tiga keputusan sekaligus -- lokasi
infrastruktur, penjadwalan kendaraan, dan manajemen pengisian daya -- ke
dalam satu kerangka optimasi, dan menunjukkan pada studi kasus Salt Lake
City bahwa penggabungan ketiganya menurunkan total biaya kepemilikan
dibandingkan skema yang mempertahankan pola rute bus diesel apa adanya.

Studi yang paling dekat dengan konteks tugas akhir ini datang dari
\textcite{sebastiani2016}, yang mengevaluasi
operasi bus listrik pada sistem \textit{Bus Rapid Transit} (BRT) dunia nyata
melalui simulasi optimasi. Karakteristik BRT -- rute tetap, target
\textit{headway} yang eksplisit (yaitu jeda waktu antara keberangkatan
atau kedatangan dua bus berurutan pada koridor yang sama), volume
penumpang tinggi -- membuat konsekuensi kegagalan
pengisian daya (bus kehabisan daya, mengantre, atau memotong rit) berdampak
langsung pada keteraturan layanan yang terukur, persis yang menjadi
perhatian utama tugas akhir ini terhadap TransJogja.

\section{Model Konsumsi Energi Baterai Bus Listrik}
\label{sec:baterai-bus-listrik}

Menjawab persoalan pada Subbab \ref{sec:bus-listrik} membutuhkan model yang
bisa memperkirakan seberapa banyak energi baterai terpakai sepanjang rute,
karena dari situlah kebutuhan pengisian daya diturunkan. Pemodelan konsumsi
energi kendaraan listrik pada literatur transportasi umumnya terbagi ke
tiga tingkat detail: mikroskopik, mesoskopik, dan makroskopik, yang berbeda
pada seberapa rinci data masukan yang dibutuhkan.

Pada ujung yang paling rinci, \textcite{basma2020} membangun
model fisis mendetail (menggunakan perangkat lunak Dymola) yang
mensimulasikan bukan hanya sistem penggerak bus, tapi juga beban termal
kabin dan sistem pendingin udara secara terpisah -- pendekatan ini akurat,
tetapi membutuhkan data operasional yang sangat rinci (profil kecepatan
detik-per-detik, kondisi cuaca, beban penumpang) yang jarang tersedia
lengkap di luar kondisi eksperimen. \textcite{chen2021} mengambil pendekatan berbeda pada tingkat detail yang
serupa: alih-alih model fisis, mereka melatih model pembelajaran mesin
(\textit{Long Short-Term Memory} dan Artificial Neural Network) dari data
pemantauan bus listrik riil di Chattanooga, Amerika Serikat, dan mencapai
galat prediksi di bawah 5\%. Kedua pendekatan ini menegaskan bahwa
akurasi tinggi pada level mikroskopik selalu berbanding lurus dengan
kebutuhan data yang sangat rinci -- baik lewat model fisis maupun model
data-driven.

Pada ujung yang lebih sederhana, model makroskopik hanya memakai kecepatan
rata-rata (bahkan kerap diasumsikan konstan) sebagai masukan utama pada
tiap segmen jalan. \textcite{montanino2025} menunjukkan bahwa
asumsi kecepatan konstan ini justru menjadi sumber bias yang signifikan --
secara analitis dan lewat eksperimen simulasi Monte Carlo berskala besar
-- yaitu menjalankan simulasi yang sama berulang kali dengan nilai acak
berbeda tiap perulangan, lalu mengagregasi hasilnya untuk mengestimasi
perilaku sistem di bawah ketidakpastian \parencite{law2015} --
mereka membuktikan konsumsi energi kendaraan listrik meningkat seiring
variansi akselerasi, sehingga model yang mengabaikan dinamika lalu lintas
(percepatan dan perlambatan) secara sistematis meremehkan konsumsi energi
riil. Mereka mengusulkan variansi kecepatan -- yang bisa diturunkan dari
kepadatan dan arus lalu lintas terukur -- sebagai proksi yang lebih baik
dibandingkan kecepatan rata-rata semata. \textcite{mamarikas2025}
mengusulkan pendekatan mesoskopik sebagai jalan
tengah antara model makroskopik yang terlalu sederhana dan model
mikroskopik yang terlalu rinci: memanfaatkan variabel lalu lintas pada
level detail menengah (jumlah dan durasi kejadian berhenti sepanjang rute)
untuk memperkirakan konsumsi energi bus listrik, tanpa membutuhkan profil
kecepatan super-detail seperti model mikroskopik, dan menunjukkan galat
prediksi yang mendekati model mikroskopik (di bawah 5\% pada data lalu
lintas kota) jauh lebih baik dibandingkan model makroskopik konvensional
(galat hingga 20\% pada kondisi tertentu). Pendekatan mesoskopik ini relevan
untuk tugas akhir ini karena data yang tersedia untuk 20 koridor TransJogja
adalah data rute (jarak antar-halte, target headway), bukan profil
kecepatan detik-per-detik tiap bus maupun data pemantauan historis yang
dibutuhkan pendekatan \textit{machine learning} -- sehingga model tingkat
mesoskopik menjadi pilihan yang paling sesuai dengan ketersediaan data.

Satu faktor yang tidak tercakup pada model mesoskopik di atas namun terbukti
berpengaruh besar pada literatur lain adalah topografi rute. \textcite{alogaili2020} mengintegrasikan model elevasi digital (\textit{Digital
	Elevation Model}) dengan model dinamika kendaraan pada studi kasus jaringan
BRT Rapid KL di Kuala Lumpur -- konteks yang cukup dekat dengan TransJogja,
karena sama-sama sistem BRT dengan rute tetap -- dan menemukan elevasi jalan
sebagai salah satu faktor utama yang membedakan konsumsi energi antar-rute.
Studi yang sama juga mengevaluasi \textit{opportunity charging} dan
\textit{overnight/depot charging} sebagai dua protokol yang dibandingkan,
sejalan dengan mekanisme yang menjadi fokus tugas akhir ini (lihat Subbab
\ref{sec:perencanaan-infrastruktur}).

Tugas akhir ini tidak memodelkan elevasi secara eksplisit -- data topografi
untuk 20 koridor TransJogja tidak tersedia dalam bentuk yang dapat
diintegrasikan pada rentang waktu penelitian ini -- dan hal ini
didokumentasikan sebagai batasan (lihat Batasan Masalah, Bab I), bukan
diabaikan begitu saja. Temuan Al-Ogaili dkk. di atas justru menjadi alasan
untuk mencatat keterbatasan ini secara eksplisit, alih-alih menganggapnya
tidak berpengaruh.

\section{Perencanaan Infrastruktur Pengisian Daya}
\label{sec:perencanaan-infrastruktur}

Dengan model konsumsi energi tersedia, pertanyaan berikutnya adalah
bagaimana menguji berbagai konfigurasi lokasi, kapasitas, dan mekanisme
pengisian daya sebelum benar-benar dibangun di lapangan. Dua pendekatan
besar muncul pada literatur untuk menjawab pertanyaan ini: optimasi
matematis dan simulasi.

\subsection{Pendekatan Optimasi Matematis}

Sebagian besar studi perencanaan infrastruktur bus listrik memakai
formulasi \textit{Mixed-Integer Linear Programming} (MILP) atau variannya
untuk menentukan lokasi \textit{charger}, ukuran armada, dan jadwal
pengisian daya secara sekaligus. \textcite{zhou2024}, dalam
tinjauan sistematis atas lebih dari 100 studi perencanaan fasilitas
\textit{charging} bus listrik, mencatat bahwa mayoritas studi memakai model
optimasi deterministik (MILP, \textit{mixed-integer nonlinear programming},
atau variannya) -- dari belasan studi representatif yang mereka bandingkan,
hanya satu-dua yang secara eksplisit memakai formulasi stokastik untuk
menangani ketidakpastian waktu tempuh atau degradasi baterai. Model
\textit{deterministik} memakai satu nilai tetap (mis. rata-rata) untuk
tiap parameter masukan, sehingga hasilnya selalu sama pada tiap
perulangan, sedangkan model \textit{stokastik} memperlakukan parameter
yang tidak pasti sebagai variabel acak berdistribusi tertentu, sehingga
perlu diulang berkali-kali untuk mendapat gambaran distribusi hasil alih-alih
nilai tunggal. Temuan ini
konsisten dengan pengamatan \textcite{guschinsky2021}, yang
secara terus terang menyatakan dalam model optimasi armada dan
infrastruktur \textit{fast-charging} mereka bahwa sebagian besar parameter
masukan bersifat tidak pasti pada kondisi riil, namun nilai deterministiknya
(nominal, rata-rata statistik, atau skenario terburuk) tetap dipakai sebagai
penyederhanaan karena kompleksitas model yang sudah tinggi.

\textcite{rizopoulos2025} merumuskan persoalan
ini sebagai model \textit{bi-objective MILP}, dan menemukan trade-off
mendasar antara meminimalkan waktu perjalanan tanpa penumpang
(\textit{deadhead time}) dan biaya pemasangan \textit{charger} pada studi
kasus Manhattan dan Limassol. \textcite{liu2021}
memperluas pertimbangan pada aspek kelistrikan -- mencocokkan daya
pengisian dengan kemampuan penerimaan baterai serta pengaruh musim (suhu
udara) terhadap performa baterai -- pada studi kasus 17 rute bus di
Beijing. \textcite{manzolli2022} mendekati
sisi lain persoalan yang sama -- interaksi baterai bus dengan jaringan
listrik (\textit{Vehicle-to-Grid}/V2G) -- dan menunjukkan skema ini bisa
ekonomis pada kondisi tertentu, meski validasinya baru diuji pada armada
kecil (11 unit bus di Coimbra, Portugal).

Satu studi yang metodologinya paling dekat dengan pendekatan tugas akhir
ini datang dari \textcite{rios2014}, yang mengusulkan model tiga
tahap untuk menghitung profil beban listrik dari sebuah stasiun pengisian
daya: tahap pertama memakai simulasi Monte Carlo untuk memodelkan
konsumsi energi harian tiap bus berdasarkan karakteristik rute, tahap kedua
memakai \textit{Discrete Event Simulation} untuk memodelkan antrean bus
pada charger yang jumlahnya terbatas, dan tahap ketiga menghitung profil
beban dari kedua hasil tersebut. Kombinasi Monte Carlo dan DES ini sejalan
dengan pendekatan yang diadopsi pada tugas akhir ini (lihat Bab IV) --
bedanya, Rios dkk. memakai kombinasi ini untuk menghitung kebutuhan
kapasitas kelistrikan jaringan (di Bogotá, Kolombia), sedangkan tugas akhir
ini memakainya untuk mengevaluasi dampak terhadap keteraturan jadwal
(headway) pada sistem BRT.

\subsection{Pendekatan Simulasi}

Pendekatan kedua memakai simulasi untuk menguji konfigurasi tanpa
menyelesaikannya sebagai persoalan optimasi matematis. \textit{Discrete
	Event Simulation} (DES) -- yang memajukan waktu dari satu kejadian ke
kejadian berikutnya, alih-alih dalam langkah waktu tetap -- banyak dipakai
karena kejadian yang relevan pada sistem bus listrik (kedatangan di SPKLU,
mulai dan selesainya sesi pengisian daya, keberangkatan rit) terjadi pada
waktu yang tidak beraturan. \textcite{jefferiesgohlich2020}
membangun perangkat simulasi berbasis DES untuk menentukan kebutuhan
armada, infrastruktur, dan energi pada jaringan bus riil, dengan
\textit{depot} dan \textit{opportunity charging} sebagai dua opsi yang
dibandingkan. \textcite{sebastiani2016} memakai pendekatan
serupa untuk mengevaluasi operasi bus listrik pada jaringan BRT,
mengonfirmasi bahwa DES cukup fleksibel untuk merepresentasikan antrean dan
kontensi sumber daya pada SPKLU yang dipakai bersama banyak rute -- persis
persoalan yang dihadapi TransJogja jika elektrifikasi diperluas ke seluruh
20 koridor sekaligus (lihat Bab III).

\textcite{casella2024} menunjukkan variasi
lain dari pendekatan DES: alih-alih memakai DES murni sebagai alat evaluasi
skenario, mereka membangun model optimasi penjadwalan pengisian daya
berbasis kejadian diskrit, dengan dua elemen yang relevan bagi tugas akhir
ini -- representasi profil pengisian daya baterai yang non-linear
(bukan asumsi linier sederhana), dan pengenalan \textit{"intermediate
	events"} yang memungkinkan sistem tidak selalu memaksa ada bus yang
sedang mengisi daya pada tiap rentang waktu, sehingga jadwal pengisian daya
punya keleluasaan waktu yang lebih realistis. Kedua elemen ini konsisten
dengan pertimbangan yang diambil pada perancangan model drain baterai
tugas akhir ini (lihat Bab IV), meski tugas akhir ini tidak mengadopsi
kerangka optimasinya secara langsung.

Satu keterbatasan yang umum ditemui pada studi simulasi di atas: kondisi
operasional yang diasumsikan bersifat deterministik atau berbasis
rata-rata, padahal waktu tempuh antar-halte pada praktiknya bervariasi dari
hari ke hari akibat kepadatan lalu lintas, cuaca, dan faktor lain yang
tidak sepenuhnya bisa diprediksi. \textcite{avishan2023} menunjukkan besarnya konsekuensi dari mengabaikan
variasi ini: pada penjadwalan armada bus listrik yang mengasumsikan waktu
tempuh tetap, hingga 48\% rit yang dijadwalkan berpotensi terlambat atau
batal begitu waktu tempuh sesungguhnya bervariasi -- angka yang jauh dari
dapat diabaikan. \textcite{jiang2021} menemukan pola
serupa pada penjadwalan bus listrik di Shenzhen, dan mengatasinya lewat
optimasi \textit{robust} yang divalidasi lewat simulasi Monte Carlo,
meski hasilnya menunjukkan ada \textit{trade-off}: keandalan jadwal yang
lebih tinggi terhadap ketidakpastian umumnya menaikkan biaya operasional.

Variasi waktu tempuh ini berpotensi memengaruhi kapan persisnya bus
membutuhkan pengisian daya, dan karenanya memengaruhi seberapa parah
kontensi yang terjadi di SPKLU pada hari-hari tertentu. \textcite{debruin2023} mengambil
pendekatan yang dekat dengan tugas akhir ini untuk menangani persoalan
serupa: alih-alih menambah kelonggaran waktu (\textit{slack}) pada jadwal
secara statis -- yang menurut mereka tidak cukup menangkap variansi
distribusi waktu tempuh maupun ketergantungan antar-rit -- mereka memakai
\textit{Discrete Event Simulation} untuk secara langsung mengevaluasi
ketahanan (\textit{robustness}) suatu jadwal terhadap waktu tempuh yang
stokastik. Lapisan simulasi stokastik Monte Carlo -- menjalankan simulasi
berulang kali dengan variasi acak pada parameter waktu tempuh, lalu
mengamati distribusi hasilnya -- menjadi pelengkap yang relevan atas
keterbatasan pendekatan deterministik di atas, dan menjadi dasar pendekatan
yang diadopsi pada tugas akhir ini (lihat Bab IV). Tugas akhir ini secara konsisten mengadopsi pendekatan simulasi \textit{stokastik}
berbasis Monte Carlo -- bukan deterministik -- karena parameter waktu tempuh
antar-halte diperlakukan sebagai variabel acak (lihat Subbab
\ref{sec:perencanaan-infrastruktur} dan Bab IV), sehingga seluruh hasil pada
tugas akhir ini dilaporkan sebagai rata-rata dan simpangan baku dari banyak
kali perulangan, bukan nilai tunggal.

\section{Posisi Tugas Akhir Ini terhadap Peta Riset yang Ada}
\label{sec:posisi-penelitian}

Tabel \ref{tab:peta-riset} merangkum tiga arus riset yang tergambar dari
subbab-subbab di atas, beserta keterbatasan yang konsisten muncul pada
masing-masing.

\begin{table}[h]
	\centering
	\caption{Tiga arus riset perencanaan bus listrik dan keterbatasannya}
	\label{tab:peta-riset}
	\begin{tabular}{p{3.5cm}p{4.5cm}p{4.5cm}}
		\toprule
		\textbf{Arus Riset} & \textbf{Capaian} & \textbf{Keterbatasan yang Konsisten Muncul} \\
		\midrule
		Optimasi infrastruktur \& penjadwalan (MILP/MINLP) &
		Matang secara metodologis untuk menentukan lokasi, ukuran armada, dan
		jadwal secara simultan \parencite{he2023,rizopoulos2025,guschinsky2021,liu2021} &
		\textcite{zhou2024} mencatat, dari tinjauan sistematis
		lebih dari 100 studi, mayoritas memakai formulasi deterministik;
		\textcite{guschinsky2021} secara terbuka mengakui
		penyederhanaan ini dalam model mereka sendiri \\
		\addlinespace
		Manajemen baterai \& interaksi jaringan (V2G) &
		Terbukti ekonomis pada kondisi tertentu \parencite{manzolli2022} &
		Validasi baru pada armada berskala kecil, belum pada armada kota
		penuh \\
		\addlinespace
		Simulasi (DES) untuk perencanaan infrastruktur &
		Mampu merepresentasikan antrean dan kontensi sumber daya pada sistem
		nyata, termasuk BRT \parencite{jefferiesgohlich2020,sebastiani2016,rios2014,casella2024} &
		Umumnya masih memakai kondisi operasional deterministik/rata-rata;
		pengecualian pada studi penjadwalan yang secara eksplisit menangani
		ketidakpastian \parencite{avishan2023,jiang2021,debruin2023}, meski
		studi tersebut berfokus pada penjadwalan kendaraan (\textit{vehicle
			scheduling}), bukan perbandingan skenario lokasi/kapasitas
		infrastruktur \\
		\bottomrule
	\end{tabular}
\end{table}

Perbandingan pada Tabel \ref{tab:peta-riset} menunjukkan bahwa perhatian
terhadap ketidakpastian operasional sesungguhnya sudah muncul pada
literatur -- terutama pada studi penjadwalan kendaraan (\textit{vehicle
	scheduling}) yang secara eksplisit memodelkan waktu tempuh sebagai variabel
acak \parencite{avishan2023,jiang2021,debruin2023}, bahkan salah satunya
\parencite{avishan2023} menunjukkan dampak yang cukup besar (hingga 48\% rit
terdampak) bila ketidakpastian ini diabaikan. Namun perhatian yang sama
belum banyak terlihat pada studi perbandingan lokasi dan kapasitas
infrastruktur SPKLU -- arus riset ini (baris pertama Tabel
\ref{tab:peta-riset}) secara konsisten tetap memakai parameter deterministik,
sebagaimana dicatat baik oleh \textcite{zhou2024} pada tinjauan
sistematisnya maupun diakui langsung oleh \textcite{guschinsky2021}
dalam model mereka sendiri.

Tugas akhir ini tidak diposisikan sebagai kontribusi pada arus optimasi
matematis (bukan MILP/MINLP, bukan pula manajemen baterai atau V2G --
keduanya di luar Batasan Masalah pada Bab I). Posisinya berada pada arus
ketiga: simulasi berbasis DES, dengan lapisan stokastik Monte Carlo yang
menjawab keterbatasan "kondisi operasional deterministik" -- bukan pada
persoalan penjadwalan kendaraan seperti studi-studi yang sudah menangani
ketidakpastian di atas, melainkan pada persoalan perbandingan lokasi dan
kapasitas infrastruktur SPKLU, yang berdasarkan penelusuran literatur pada
bab ini belum banyak diuji dengan cara yang sama. Preseden metodologis yang paling dekat adalah \textcite{rios2014}, yang memakai kombinasi
serupa (Monte Carlo untuk konsumsi energi harian, DES untuk antrean SPKLU)
-- namun untuk tujuan menghitung profil beban jaringan listrik, bukan
mengevaluasi keteraturan jadwal (headway) seperti tugas akhir ini.

Tugas akhir ini juga tidak mengklaim menyelesaikan seluruh keterbatasan
yang ada pada literatur -- kombinasi lokasi/kapasitas SPKLU pada tugas
akhir ini tetap diuji lewat perbandingan skenario (bukan dioptimasi
otomatis lewat MILP), sebagaimana disebutkan pada Batasan Masalah. Namun
pada aspek representasi variasi harian untuk persoalan perencanaan
infrastruktur secara spesifik, tugas akhir ini melangkah lebih jauh
dibandingkan pendekatan deterministik yang selama ini umum dipakai pada
arus riset tersebut.

Dari sisi konteks, studi kasus pada literatur yang ditinjau di atas
sebagian besar berasal dari Eropa, Amerika Serikat, Tiongkok, dan sebagian
Asia Tenggara (Malaysia); sistem BRT di Indonesia -- termasuk TransJogja --
belum banyak menjadi objek kajian serupa. Tugas akhir ini turut mengisi
celah kontekstual tersebut, meski hal ini bersifat kontribusi
empiris/kontekstual, bukan kebaruan metodologis.